% Junior Level 1, Buffer session B2 -- Remix Studio
\kcpart{B2}{Remix Studio}{Take something you made, change it on purpose, and say what you changed.}
\lessoninfo{Concepts: remixing (I), all concepts so far (R), giving credit (R)}{Tools: Scratch 3 offline app, a dice, saved projects from earlier weeks}{CSTA: 1A-AP-10, 1A-AP-13, 1A-AP-14, 1A-AP-15}

\begin{teachernote}[When to use B2]
Use B2 any time after week 5: after a missed session, before a holiday, or between weeks 9 and 10 as extra studio time to finish greeting cards. It needs at least one saved project from an earlier week (for example Hello Cat, a week 4 animation or a week 8 drawing). If nothing was saved, rebuild Hello Cat from week 2 in the first five minutes of Make it.
\end{teachernote}

\begin{objectives}
  \item The student can choose a past project and change it in two ways on purpose. (J1.G2) [J1.B2.1]
  \item The student can add a click or key event to an old project so it does something new. (J1.G3) [J1.B2.2]
  \item The student can find and fix a bug that a partner hid in their remix. (J1.G4) [J1.B2.3]
  \item The student can say what they changed, and give credit when the idea came from someone else. (J1.G6, J1.G7) [J1.B2.4]
\end{objectives}

\begin{materials}
  \item One dice (or six folded slips of paper in a cup, numbered 1 to 6). The remix menu is in the box \enquote{Remix dice} below.
  \item On each Mac: the Scratch app open, and the children's saved projects together in one folder (for example \enquote{Documents/scratch}). A mouse plugged in.
  \item A sand timer (5 minutes). Prep time: about 5 minutes.
\end{materials}

\needspace{16\baselineskip}
\section{At a glance}
\begin{timeplan}
  \timeslot{0--8}{Warm-up}{Remix a dance: do the week 4 dance, change two moves, and the others guess what changed. Then roll the remix dice.}
  \timeslot{8--14}{Show me}{Open Hello Cat, save it with a new name, change two things, and say what changed and whose idea it was.}
  \timeslot{14--36}{Make it}{Pairs remix a past project with two dice twists (14--28), hidden bug swap (28--32), free remix (32--36).}
  \timeslot{36--38}{Wiggle break}{Remix the wiggle: everyone adds one new move to a stretch.}
  \timeslot{38--45}{Share}{\enquote{I started from \ldots. I changed \ldots{} and \ldots.} Credit check, exit question, reflection.}
\end{timeplan}

\section{Warm-up (0--8 min)}
\begin{unplugged}[Remix a dance]
\begin{enumerate}[leftmargin=1.4em,itemsep=1pt]
  \item Do the dance from week 4 together once (for example: repeat 4 of clap, stomp, spin).
  \item One child changes two things: a different move, or a different number in the repeat. They do the new dance.
  \item The others guess: \enquote{What did you change?} \enquote{You did repeat 2, and you jumped instead of stomping!}
  \item \enquote{That is a \textbf{remix}: we start from something that already exists, and we change it on purpose. And we say where it came from: \enquote{This is a remix of our week 4 dance.}}
\end{enumerate}
With 1 child: you remix first and the child guesses, then swap. With 2--4: everyone remixes once, in turn. Then roll the dice once each to see the first twist for today.
\end{unplugged}

\begin{tip}[Remix dice]
Roll the dice and add that twist to your project.
\textbf{1:} a new sound (pink Sound). \textbf{2:} a new backdrop (purple Looks, \enquote{switch backdrop to}). \textbf{3:} a dance with a repeat (orange Control). \textbf{4:} a click or key event (yellow Events). \textbf{5:} draw or stamp with the pen (green Pen). \textbf{6:} your choice!
\end{tip}

\section{Show me (8--14 min)}
\begin{teacherscript}
\enquote{First, we keep the old project safe.} Click \menu{File} → \menu{Load from your computer} and open Hello Cat. Then click \menu{File} → \menu{Save to your computer} and save it with a new name: \enquote{hello-cat-remix}. \enquote{Now the old one is safe, and we can change this one as much as we like.}
\end{teacherscript}

\begin{analogy}[A new recipe from an old one]
\enquote{A remix is like taking a pancake recipe and adding blueberries. It is still pancakes, but now it is your version. When you share it you say: \enquote{It is Grandma's recipe, with my blueberries.}}
\end{analogy}

\begin{workedexample}[Remix Hello Cat with two twists]
Roll the dice twice. Say you get 3 (a dance) and 4 (a click event).
\begin{enumerate}[leftmargin=1.4em,itemsep=1pt]
  \item Twist 3: click the orange \menu{Control} circle. Snap \enquote{repeat 10} under the cat's script and change 10 to 4.
  \item Click the purple \menu{Looks} circle. Put \enquote{next costume} inside the repeat. Click the orange \menu{Control} circle and put \enquote{wait 1 seconds} under it.
  \item Twist 4: click the yellow \menu{Events} circle. Drag \enquote{when this sprite clicked} to an empty space. Click the pink \menu{Sound} circle and snap \enquote{start sound Meow} under it.
  \item Click the green flag: the cat says hi, meows, moves and dances. Click the cat: it meows.
  \item Say it out loud: \enquote{I started from my Hello Cat. I added a dance and a click.} Save again.
\end{enumerate}
\begin{center}
\begin{scratch}[scale=0.85]
\blockinit{when \greenflag clicked}
\blocklook{say \ovalnum{Hi!} for \ovalnum{2} seconds}
\blocksound{play sound \selectmenu{Meow} until done}
\blockmove{move \ovalnum{50} steps}
\blockrepeat{repeat \ovalnum{4}}{
  \blocklook{next costume}
  \blockcontrol{wait \ovalnum{1} seconds}
}
\end{scratch}
\quad
\begin{scratch}[scale=0.85]
\blockinit{when this sprite clicked}
\blocksound{start sound \selectmenu{Meow}}
\end{scratch}
\end{center}
\end{workedexample}

\section{Make it (14--36 min)}
\begin{pairroles}[Driver and Helper]
The \textbf{Driver} holds the mouse and makes the change. The \textbf{Helper} rolls the dice, reads the twist from the remix menu and points to the right color circle, but never grabs the mouse. Swap on the timer (5 minutes). With 1 child: you are the Helper and roll the dice. With 3: Driver, Helper and Checker (clicks the flag to test each twist). With 4: two pairs, one project per pair, then swap projects for the hidden bug swap.
\end{pairroles}

Each pair picks a past project, saves it with a new name, and adds two dice twists. Test after each twist.

\begin{bughunt}[Hidden bug swap (minutes 28--32)]
Each owner saves their remix first. Then the partner asks, \enquote{May I hide a bug?} and makes one small change while the owner looks away: pull one block off a script, move a block out of a repeat, or change the key in \enquote{when key pressed}. The owner clicks the green flag, tests the click or key, and hunts: \enquote{What should happen? What happened? Which block is the clue?} Then they fix it. With 1 child, you hide the bug.
\end{bughunt}

\begin{freebuild}[Free remix]
Last 4 minutes: roll once more, or pick any twist you like. Remixing a greeting card is fine too (add a pen border, a new sprite or a key that plays a sound). Save at the end.
\end{freebuild}

\section{Wiggle break (36--38 min)}
Everyone stands up and does a stretch together: reach up, touch toes, twist. Then each child remixes it by adding one new move, and everyone copies.

\section{Share (38--45 min)}
\begin{teacherscript}
Each child clicks the green flag and shows their remix. They use the sentence: \enquote{I started from \ldots. I changed \ldots{} and \ldots.} If the idea came from a partner: \enquote{This is a remix of Zoe's dancing cat.} Others say which change they liked.
\end{teacherscript}

\begin{safety}[Credit and asking first]
\enquote{When we remix our own project, we say what we started from. When we remix someone else's idea, we give them credit: we say their nickname. And we only change someone else's project if they say yes, on a copy with a new name. We remix projects from our family, not things from strangers on the internet.} Ask: \enquote{Kai wants to remix Mei's card. What does he do first?} (Ask Mei, save a copy with a new name, and say \enquote{This is a remix of Mei's card} when he shows it.)
\end{safety}

\begin{formativecheck}[Exit question]
\textbf{Look for:} the project was saved with a new name before changing it; two changes that work; the child can name both changes; credit given when the idea came from someone else.

\textbf{Ask:} \enquote{Why do we save the remix with a new name?}
\textbf{Answer:} so the old project stays safe and does not change.
\end{formativecheck}

\begin{reflection}
Ask: \enquote{Which twist was the most fun?} and \enquote{Whose idea would you like to remix next time?} Children give a thumbs up, sideways or down for how the remix went.
\end{reflection}

\section{Notes for the teacher}
\begin{commonmistakes}
  \item \textbf{Changing the old project.} Always \menu{File} → \menu{Save to your computer} with a new name before changing anything.
  \item \textbf{Too many changes at once.} Test after each twist, so a bug is quick to find.
  \item \textbf{A new script with no hat block.} The click or key twist needs \enquote{when this sprite clicked} or \enquote{when key pressed} on top.
  \item \textbf{No Pen blocks for twist 5.} Add the Pen extension (\menu{Add Extension}, bottom-left).
  \item \textbf{Forgetting credit.} Make the sentence \enquote{I started from \ldots} part of every share.
\end{commonmistakes}

\begin{differentiation}
\supportpath Use Hello Cat and make one change together, with you as the Helper. Roll the dice but allow a re-roll if the twist feels too big.
\extensionpath Make three twists, or remix a partner's project (with permission) and give credit in a \enquote{say} block at the start: \enquote{Remix of Leo's cat}.
\end{differentiation}

\begin{stretchchallenge}[Remix back]
Swap remixes with a partner (on copies). Each of you adds one twist to the other's remix, then show both versions side by side.
\end{stretchchallenge}

\begin{teachernote}[If you are short on time]
Skip the dance remix and roll the dice straight away. One twist instead of two. Keep the share sentence with credit: it is the heart of a remix.
\end{teachernote}
